\section{Devotion to the Angels and Its Regulation}\label{sec:devotion}

The faithful carry into their own prayer what the Church prays at the
altar. They greet their guardian angel morning and evening, call on
Michael against the enemy, count the nine choirs on a chaplet, and go on
pilgrimage to the mountains where the Archangel is honoured. The Church
has fostered these devotions, enriching them with indulgences, and she has
ruled them from the beginning, holding the cult of the angels to what
Scripture gives and to the worship of the one God. Her Directory of 2001
states both halves of that practice: the Church, which ``constantly
experiences their `mysterious and powerful assistance,'\,'' venerates the
angels and ``has recourse to their prompt intercession,'' and ``down
through the centuries, the faithful have translated into various
devotional exercises the teaching of the faith in relation to the
ministry of Angels'' (\emph{Directory on Popular Piety} 215--216). The
doctrine on which the devotion rests is the Church's, restated in the
Directory in the Catechism's words: ``the existence of the spiritual,
non-corporeal beings that Sacred Scripture usually calls `angels' is a
truth of faith'' (\emph{Directory} 213; CCC~328; \S\ref{sec:faith}; for the acts and their dates,
Appendix~\ref{app:chronology}).\footnote{The prayers are
cited from the \emph{Raccolta}, the Roman collection of indulgenced
prayers, in Ambrose St John's English (London: Burns and Oates, 1910),
which records each grant with its pontiff and date; from the Vatican
\emph{Collectio precum piorumque operum} of 1929, which gives the Latin of
the prayers after Low Mass; and from the American \emph{Key of Heaven}
(Baltimore: Murphy, 1901), which gives their English.}

\sectionguard
\subsection{The Prayer to the Guardian Angel}

The Church's prayer to the guardian angel is the \latin{Angele Dei}.
The \emph{Raccolta} prints it in Latin with a close English version, and
adds a rhymed English version:
\begin{quote}
\latin{Angele Dei, qui custos es mei, me tibi commissum pietate superna
illumina, custodi, rege, et guberna. Amen.}

O Angel of God, whom God hath appointed to be my guardian, enlighten and
protect, direct and govern me. Amen.

Angel of God, my guardian dear,\\
To whom his love commits me here,\\
Ever this day be at my side,\\
To light and guard, to rule and guide. Amen.
\end{quote}
(\emph{Raccolta}, 1910, n.~296, p.~267.) The \emph{Raccolta} records the
indulgences granted to it: a hundred days each time; a plenary indulgence
on 2 October, the feast of the Guardian Angels, for those who have said
it morning and evening through the year; a plenary indulgence once a
month; and a plenary indulgence at the hour of death for those who have
said it often. The grants it cites are those of Pius~VI (briefs of
2 October 1795 and 11 June 1796) and Pius~VII (15 May 1821; ibid.,
pp.~266--267). The Directory of 2001 calls the \latin{Angele Dei}
``especially popular,'' a prayer ``often recited by families at morning
and evening prayers, or at the recitation of the Angelus'' (\emph{Directory}
216).

The prayer is a compact treatise on the guardian's office. It confesses
first that the one who prays has been entrusted to the angel, \latin{tibi
commissum}, and that the trust is an act of heavenly love, \latin{pietate
superna}: the guardianship comes from the ``unspeakable providence'' of
the collect of 2 October, which sends the holy angels \latin{ad nostram
custodiam} (\emph{BR}, 2 Oct.; \S\ref{sec:liturgy}). It then asks
four things. \latin{Illumina}: the angel enlightens, for Aquinas teaches
that the angels strengthen the human intellect and propose intelligible
truth to it ``under the similitudes of sensible things'' (\ST{I q.~111
a.~1 co.}; \S\ref{sec:government}). \latin{Custodi}: the angel guards, for
every man from his birth has ``an Angel appointed to guard it,'' as Jerome
teaches in the lesson the Church reads on 29 September (\emph{In Matth.}
III, on 18:10; \ST{I q.~113 a.~5 s.c.}; \S\ref{sec:guardians}).
\latin{Rege} and \latin{guberna}: the angel rules and governs, the verbs of
the providential order in which, as Aquinas puts it, ``the order of
Divine Providence disposes that lower things be subject to the actions
of higher'' (\ST{I q.~111 a.~1 co.}; \ST{I q.~110 a.~1}). The rhymed English renders the four
verbs ``To light and guard, to rule and guide.''

Bernard gave the Church the attitude with which the prayer is to be said.
The lesson the Roman Breviary reads from his sermon on the ninetieth psalm
asks what ``this word,'' \latin{Angelis suis mandavit de te}, ought to
work in the Christian: ``What respect, what thankfulness, what trust,
ought this word to work in thee! Respect for their presence, thankfulness
for their kindness, trust in their safe-keeping'' (Bute, IV, p.~616;
\emph{BR}, 2 Oct., p.~400; \S\ref{sec:liturgy}). The Directory of
2001 draws the same three dispositions from the devotion: ``devout
gratitude to God for having placed these heavenly spirits of great
sanctity and dignity at the service of man''; ``an attitude of devotion
deriving from the knowledge of living constantly in the presence of the
Holy Angels of God''; and ``serenity and confidence in facing difficult
situations, since the Lord guides and protects the faithful in the way of
justice through the ministry of His Holy Angels'' (\emph{Directory} 216).
The Directory names Basil and Bernard as masters of the devotion. For
Basil, in the Directory's rendering, ``each and every member of the
faithful has a Guardian Angel to protect, guard and guide them through
life''; Bernard was ``a great master and a notable promoter of devotion to
the Guardian Angels'' (ibid.). The Catechism gives Basil's sentence in
its own English: ``Beside each believer stands an angel as protector and
shepherd leading him to life'' (CCC~336).

The \emph{Raccolta} adds a novena to the guardian angel, which ``may be
made at any time and with any form of prayer sanctioned by competent
ecclesiastical authority,'' with an indulgence of three hundred days each
day and a plenary indulgence once during the novena, granted by Pius~IX
on 26 November 1876 (\emph{Raccolta}, 1910, n.~297, p.~267). The Church
does not fix the words; she asks only that they be approved.

\sectionguard
\subsection{The Prayer to Saint Michael}

\paragraph{The prayers after Low Mass.} Leo~XIII ordered that certain
prayers be said after every Low Mass, kneeling, in all the churches of the
world. The \emph{Key of Heaven} of 1901 prints them under that heading:
three Hail Marys, the \latin{Salve Regina} with its versicle, the prayer
``O God, our refuge and our strength,'' asking the intercession of the
Mother of God, Saint Joseph, the Apostles, and all the saints ``for the
conversion of sinners, and for the freedom and exaltation of Holy Mother
the Church,'' and the prayer to the Archangel (\emph{Key of Heaven}, 1901,
pp.~138--141). The Vatican \emph{Collectio} of 1929 gives the Latin under
the title \latin{Preces post privatae Missae celebrationem recitandae}:
\begin{quote}
\latin{Sancte Michael Archangele, defende nos in praelio, contra nequitias
et insidias diaboli esto praesidium: Imperet illi Deus, supplices
deprecamur, tuque, Princeps militiae caelestis, Satanam aliosque spiritus
malignos, qui ad perditionem animarum pervagantur in mundo, divina
virtute in infernum detrude. Amen.} (\emph{Collectio precum}, 1929,
n.~331, p.~328)
\end{quote}
\begin{quote}
Holy Michael, the Archangel, defend us in the battle; be our protection
against the wickedness and snares of the devil---Rebuke him, O God, we
suppliantly beseech Thee: and do thou, O Prince of the heavenly host, by
the divine power drive into hell Satan and the other evil spirits, who
wander through the world seeking the ruin of souls. Amen. (\emph{Key of
Heaven}, 1901, pp.~140--141)
\end{quote}
The \emph{Collectio} records the indulgence of three hundred days and
cites the decrees of the Congregation of Rites of 6 January 1884 and 24
November 1915; the \emph{Key of Heaven} states that Leo~XIII ``grants to
all who recite these prayers, as aforesaid, 300 days indulgence''
(\emph{Collectio}, p.~328; \emph{Key of Heaven}, p.~141).

Almost every phrase of the prayer is Scripture or the liturgy. \latin{Defende
nos in praelio} is the Alleluia of Michaelmas, \latin{Sancte Michael
Archangele, defende nos in praelio: ut non pereamus in tremendo judicio}
(\emph{MR}~1920, 29 Sept.; \S\ref{sec:liturgy}), which the
\emph{Raccolta} also gives as an indulgenced antiphon: ``St Michael
Archangel, defend us in the day of battle, that we may not be lost in the
dreadful judgment'' (\emph{Raccolta}, 1910, n.~293, p.~266; Leo~XIII,
19 August 1893). The battle is the Apostle's: ``Put you on the armour of
God, that you may be able to stand against the deceits of the devil. For
our wrestling is not against flesh and blood; but against principalities
and powers, against the rulers of the world of this darkness, against the
spirits of wickedness in the high places'' (Eph 6:11--12). \latin{Imperet
illi Deus} is the word of Michael himself in the Epistle of Jude, who
``disputing with the devil, contended about the body of Moses, he durst
not bring against him the judgment of railing speech, but said: The Lord
command thee'' (Jude 9). The Archangel does not command the devil in his
own name; he calls on God's command, and the Church prays as he prayed.
\latin{Princeps militiae caelestis} is the title of the Breviary,
\latin{princeps militiae Angelorum} (\emph{BR}, 29 Sept., fourth
responsory). The spirits ``who wander through the world seeking the ruin
of souls'' are the adversary of whom Peter writes: ``your adversary the
devil, as a roaring lion, goeth about seeking whom he may devour'' (1~Pet
5:8). And the last petition asks the Archangel to do by the power of God,
\latin{divina virtute}, what the Apocalypse shows him doing: ``that great
dragon was cast out, that old serpent, who is called the devil and Satan,
who seduceth the whole world. And he was cast unto the earth: and his
angels were thrown down with him'' (Apoc 12:9).

The theology of the prayer is Gregory's. The name of Michael, \latin{Quis
ut Deus}, answers the boast of the enemy who said \latin{similis ero
Altissimo} (Isa 14:14), and at the end of the world the old enemy ``is
shown unto us a-fighting with Michael the Archangel'' so that, as Gregory
concludes, \latin{qui se ad Dei similitudinem superbus extulerat, per
Michaelem peremptus discat, quia ad Dei similitudinem per superbiam nullus
exsurgat}: he who proudly raised himself to God's likeness is to learn,
slain by Michael, that no one rises to God's likeness by pride (Gregory,
\emph{Hom.\ in Ev.} 34.9; Bute, IV, p.~596). The prayer asks the
Archangel to fight now the battle he will finish at the end, and it asks
that the power be God's. It does not make the world a contest of equal
forces; it asks the prince of the heavenly host to act by the divine
power against creatures who are subject to it (\ST{I q.~114};
\S\ref{sec:assaults}; on the limits of the demons' power,
Appendix~\ref{app:demons}).

\paragraph{The longer prayer of Leo~XIII.} The \emph{Raccolta} also
prints a longer prayer to Saint Michael, with an indulgence of three
hundred days, citing a motu proprio of Leo~XIII of 25 September 1888. It
begins: ``O glorious Archangel St Michael, Prince of the heavenly host, be
our defence in the terrible warfare which we carry on against
principalities and Powers, against the rulers of this world of darkness,
spirits of evil,'' and ends with versicles and a prayer asking God, by the
intercession of the Virgin Mary and of the Archangel, ``to help us against
Satan and all other unclean spirits, who wander about the world for the
injury of the human race and the ruin of souls'' (\emph{Raccolta}, 1910,
n.~292, pp.~264--265).

\paragraph{The Office prayed at home.} Pius~VII enriched with an
indulgence the private recitation of the Office hymn of Michaelmas, its
Magnificat antiphon, and its collect: \latin{Te splendor et virtus
Patris}, \latin{Princeps gloriosissime, Michael Archangele, esto memor
nostri}, and \latin{Deus, qui miro ordine} (\emph{Raccolta}, 1910, n.~289,
pp.~259--260; Pius~VII, 6 May 1817). The \emph{Raccolta} prints the
antiphon in English, ``Most glorious Prince, Michael the Archangel, be
mindful of us: pray for us always both here and everywhere to the Son of
God,'' and the collect, ``O God, who disposest the services of Angels and
men in a wonderful order; mercifully grant that those who ever stand
before Thee, ministering to Thee in Heaven, may themselves also protect
our life here upon earth'' (ibid., p.~260). Private devotion here is the
liturgy carried home: the layman prays in his room what the Church sings
in choir on the Archangel's feast. The \emph{Raccolta} adds a novena of
Saint Michael, in any approved form, granted by Pius~IX on 26 November
1876 (n.~290, p.~261).

\sectionguard
\subsection{The Chaplet of Saint Michael}

The \emph{Raccolta} gives the chaplet of Saint Michael under the title
``Angelical Crown,'' with the indulgences granted by Pius~IX on 8 August
1851: seven years and seven quarantines; a hundred days daily to anyone
who carries the chaplet or kisses the medal of the holy angels attached
to it; a plenary indulgence once a month; and plenary indulgences on the
feasts of the Apparition of Saint Michael (8 May), his Dedication (29
September), Saint Gabriel (``March 18,'' as that edition prints the day),
Saint Raphael (24 October), and the Guardian Angels (2 October)
(\emph{Raccolta}, 1910, n.~291, p.~261). The chaplet opens with the
versicle of the Hours, ``O God, come to my assistance,'' and the
\latin{Gloria Patri}. It then makes nine salutations, one to each choir,
each of one \latin{Pater} and three \latin{Aves}, with a petition ``at the
intercession of St Michael'' and of the choir:
\begingroup\small
\begin{longtable}{>{\bfseries\raggedright\arraybackslash}p{0.14\linewidth}>{\raggedright\arraybackslash}p{0.20\linewidth}>{\raggedright\arraybackslash}p{0.58\linewidth}}
\toprule
\textbf{Salutation} & \textbf{Choir} & \textbf{Grace asked (\emph{Raccolta} 1910, pp.~261--263)}\\
\midrule
\endhead
First & Seraphim & ``to make us worthy to receive into our hearts the
fire of his perfect charity''\\
Second & Cherubim & ``grace to abandon the ways of sin, and follow the
path of Christian perfection''\\
Third & Thrones & ``to infuse into our hearts a true and earnest spirit of
humility''\\
Fourth & Dominations & ``grace to have dominion over our senses, and to
correct our depraved passions''\\
Fifth & Powers & ``to keep our souls from the wiles and temptations of
the devil''\\
Sixth & Virtues & ``keep us from falling into temptations and deliver us
from evil''\\
Seventh & Principalities & ``to fill our souls with the spirit of true
and hearty obedience''\\
Eighth & Archangels & ``the gift of perseverance in the faith and in all
good works, that we may thereby be enabled to attain unto the glory of
Paradise''\\
Ninth & Angels & ``that the Holy Angels may protect us during life, and
after death may lead us into the everlasting glory of heaven''\\
\bottomrule
\end{longtable}
\endgroup
Four \latin{Paters} follow, ``the first to St Michael, the second to St
Gabriel, the third to St Raphael, the fourth to our Angel Guardian,'' and
the chaplet ends with an antiphon to the Archangel, ``Michael, glorious
Prince, chief and champion of the heavenly host, guardian of the souls of
men, conqueror of the rebel angels, who art set over the palace of God,
our worthy captain under Jesus Christ,'' the versicle ``Pray for us, most
blessed Michael, Prince of the Church of Jesus Christ,'' and a collect
that asks ``that at the hour of our death no one of them may approach to
harm us, and that by the same Archangel Michael we may be introduced into
the presence of thy high and heavenly Majesty'' (\emph{Raccolta}, 1910,
p.~263).

The chaplet descends the hierarchy from the Seraphim to the Angels. Its
order agrees with Dionysius and Gregory in the first four places and the
last two, and differs from both in the middle: it sets the Powers fifth,
the Virtues sixth, and the Principalities seventh, where Dionysius has
Virtues, Powers, Principalities and Gregory has Principalities, Powers,
Virtues (Appendix~\ref{app:orders}; \ST{I q.~108 a.~6}). The graces it asks
follow the Fathers' account of the choirs. Gregory teaches that the
Seraphim are named ``burning,'' \latin{ardentes vel incendentes}, because
they burn with an incomparable love, and the first salutation asks the fire
of perfect charity. He teaches that the Powers are those to whose rule the
adverse powers are subject, \latin{quorum potestate refrenantur, ne corda
hominum tantum tentare praevaleant quantum volunt}, restrained so that
they cannot tempt the hearts of men as much as they would; the fifth
salutation asks the Powers' help against the wiles of the devil. He
teaches that the Dominations are named from ruling, and the fourth
salutation asks dominion over the senses (Gregory, \emph{Hom.\ in Ev.}
34.10).

Gregory went further, and the chaplet prays his further teaching. The
heavenly city is made of angels and men, and the ways of life of
individual men answer to the orders of the hosts, so that men are counted
into the lot of an order by likeness of life: \latin{distinctae namque
conversationes hominum singulorum agminum ordinibus congruunt, et in
eorum sortem per conversationis similitudinem deputantur}. Those who
rule their vices and desires are counted among the Dominations; those who
cast out evil spirits by the power of prayer, among the Powers; those
full of the love of God and neighbour, among the Cherubim, since the name
Cherubim means fullness of knowledge and charity is the fullness of the
law; and those who burn with desire for their Maker alone, \latin{quorum
cor in igne conversum lucet et urit}, among the Seraphim (\emph{Hom.\ in Ev.}\ 34.11). The chaplet asks, choir by choir, for the graces by
which men come to share the lot of the angels.

\sectionguard
\subsection{Gabriel and Raphael}

\paragraph{The Angelus.} The faithful remember Gabriel every day, morning,
noon, and evening, at the sound of the bell. The \latin{Angelus} opens
with his message:
\begin{quote}
\latin{V.~Angelus Domini nuntiavit Mariae. R.~Et concepit de Spiritu
Sancto. Ave Maria.}

V.~The angel of the Lord declared unto Mary. R.~And she conceived of
the Holy Ghost. Ave Maria. (\emph{Raccolta}, 1910, n.~182, p.~151)
\end{quote}
Its collect asks \latin{ut qui, angelo nuntiante, Christi Filii tui
Incarnationem cognovimus, per Passionem ejus et Crucem ad resurrectionis
gloriam perducamur}, ``that we, to whom the Incarnation of Christ thy Son
was made known by the message of an angel,'' may be brought by his
Passion and Cross to the glory of the resurrection (ibid., p.~152). The
\emph{Raccolta} records the indulgences granted to those who say it at the
bell, from Benedict~XIII's brief of 14 September 1724 to Leo~XIII (ibid.,
p.~151). The devotion is Marian, and it is the angel's: the Church greets
the Mother of God three times a day in the words the Archangel brought
her, and her collect keeps the phrase of the Mass of Saint Gabriel,
\latin{Gabriele nuntiante}, the Incarnation known by the angel's message
(\emph{MR}~1920, 24 Mar.; \S\ref{sec:liturgy}). The Directory of
2001 joins the two devotions: the \latin{Angele Dei} ``is often recited by
families~\dots\ at the recitation of the Angelus'' (\emph{Directory} 216).

\paragraph{The prayer to Saint Raphael.} The \emph{Raccolta} gives a
novena to Gabriel or Raphael in any approved form (n.~294, Pius~IX,
26 November 1876) and a prayer to Saint Raphael enriched by Leo~XIII on
21 June 1890:
\begin{quote}
O glorious Archangel, St Raphael, great Prince of the heavenly court,
illustrious for thy gifts of wisdom and grace, guide of those who journey
by land or sea, consoler of the afflicted, and refuge of sinners; I beg
thee to assist me in all my needs and in all the sufferings of this life,
as once thou didst help the young Tobias on his travels. And because thou
art the medicine of God, I humbly pray thee to heal the many infirmities
of my soul, and the ills which afflict my body, if it be for my greater
good. I specially ask of thee an angelic purity which may fit me to be
the temple of the Holy Spirit. Amen. (\emph{Raccolta}, 1910, n.~295,
p.~266)
\end{quote}
Every title of the prayer is drawn from the book of Tobias and the
Fathers' reading of it. Raphael is the guide of travellers because he went
with Tobias on his journey, ``a beautiful young man,
standing girded, and as it were ready to walk'' (Tob 5:5); he is the
consoler of the afflicted because ``the holy angel of the Lord, Raphael
was sent to heal them both, whose prayers at one time were rehearsed in
the sight of the Lord'' (Tob 3:25), the blind Tobit and the afflicted
Sara; he is ``the medicine of God'' by the interpretation of his name that
the Church reads from Gregory on Michaelmas, \latin{Raphael vero dicitur
medicina Dei}, given to him who touched as a physician the eyes of Tobias
(\emph{Hom.\ in Ev.} 34.9; Appendix~\ref{app:names}). The prayer asks
him for health of soul first and of body only ``if it be for my greater
good,'' and it ends with purity, as the book ends the danger of Sara's
marriage when ``the angel Raphael took the devil, and bound him'' (Tob
8:3): the angel who bound the demon Asmodeus (Tob 3:8) is asked to keep
the body a temple of the Holy Spirit.

\sectionguard
\subsection{The Sanctuaries of the Archangel}

The Directory of 2001 names three ``great shrines'' in honour of the
angels, ``Mont-Saint-Michel in Normandy, San Michele della Chiusa in
Piemonte and San Michele Gargano in Apulia, each appointed with specific
feast days'' (\emph{Directory} 216). The Church relates the origins of
the Gargano shrine and of Rome's sanctuary of the Archangel in her Office and
Martyrology; the medieval \emph{Golden Legend} of Jacobus de Voragine
relates both, with the apparition at Mont-Saint-Michel. The Church keeps them as pious
tradition.

\paragraph{Monte Gargano.} The Breviary reads the story of the mountain
in the second nocturn of 8 May. It begins with the principle:
\begin{quote}
\latin{Beatum Michaelem Archangelum saepius hominibus apparuisse, et
sacrorum librorum auctoritate, et veteri sanctorum traditione
comprobatur. Quamobrem multis in locis facti memoria celebratur. Eum ut
olim synagoga Judaeorum, sic nunc custodem et patronum, Dei veneratur
Ecclesia. Gelasio autem primo Pontifice Maximo, in Apulia in vertice
Gargani montis, ad cujus radices incolunt Sipontini, Archangeli Michaelis
fuit illustris apparitio.} (\emph{BR}, \emph{Pars verna}, 8 May, lesson
iv, p.~571)
\end{quote}
``That the blessed Archangel Michael hath oftentimes been seen of men is
attested on the authority of the Holy Bible, and also by the ancient
traditions of the Saints. For this reason such visions are held in
remembrance in many places. As of old time did the Synagogue of the Jews,
so now doth the Church of God venerate Michael as her watcher and
defender. But during the Popedom of Gelasius I.\ the summit of Mount
Gargano in Apulia, at whose foot lieth the town of Siponto, was the scene
of an extraordinary appearance of this same Archangel Michael'' (Bute,
II, p.~869). The Church venerates Michael as the Synagogue did: the
prince of Israel in Daniel, ``Michael your prince'' (Dan 10:21), is the
guardian of the Church.

The fifth and sixth lessons tell the sign. A bull from the herds of
Gargano strayed and was found at the mouth of a cave; one of the herdsmen
shot an arrow at it, and ``the arrow turned round and came back against
him that had shot it'' (Bute, II, p.~870; \latin{retorta sagitta in
ipsum recidit sagittarium}, \emph{BR}, p.~572). No one dared approach
the cave. The people consulted the bishop of Siponto, who proclaimed a
fast of three days and answered that the meaning was to be sought from
God. ``After three days the Archangel Michael gave warning to the Bishop
that that place was under his protection, and that he had thus pointed
out by a sign that he wished that worship should be offered to God there,
with remembrance of himself and of the Angels'' (Bute, II, p.~870). The
Latin is precise about the object of the worship: \latin{velle ibi cultum
Deo in sui et Angelorum memoriam adhiberi}, worship offered to God, in
memory of the Archangel and the angels (\emph{BR}, p.~572). The bishop and
the citizens found the cave shaped like a church and began to celebrate
the divine offices there, ``which sanctuary hath been glorified with many
miracles'' (Bute, II, p.~870). The Martyrology of 29 September remembers
the same church, ``mean as to building, but filled with power from
heaven'' (Bute, IV, p.~592).

The \emph{Golden Legend}, in Caxton's English, tells the story more
fully, dating it to the year 390 and naming the rich owner of the herds
Garganus; in its telling Michael says to the bishop, ``Know ye that this
man is so hurt by my will. I am Michael the archangel, which will that
this place be worshipped in earth, and will have it surely kept''
(\emph{Golden Legend}, ed.\ Ellis, 1900, V, pp.~182--183). The shrine, like
the feast, is a school of the doctrine the collect of Michaelmas prays:
the angel asks that God be worshipped, and keeps the place where he is.

\paragraph{Rome and the Castle of the Angel.} The same lessons of 8 May
relate that soon afterwards Pope Boniface dedicated a church of Saint
Michael at Rome on the third day before the Kalends of October,
\latin{Bonifacius Papa Romae in summo circo sancti Michaelis ecclesiam
dedicavit tertio Kalendas Octobris: quo die etiam omnium Angelorum
memoriam Ecclesia celebrat} (\emph{BR}, 8 May, lesson vi, p.~572). Bute's
English places the church ``on Hadrian's Mole'' (Bute, II, p.~870).
Hadrian's mausoleum took its Christian name from the Archangel in the
tradition the \emph{Golden Legend} records under his feast. In the time
of Gregory the Great, when the pope ``had established the litanies for the
pestilence that was that time, and prayed devoutly for the people, he saw
upon the castle which was said sometime: The memory of Adrian, the angel
of God, which wiped and made clean a bloody sword, and put it into a
sheath. And thereby he understood that his prayers were heard. Then he
did do make there a church in the honour of S.\ Michael, and that castle
is yet named the Castle Angel'' (\emph{Golden Legend}, V, p.~184). The
Castel Sant'Angelo keeps the memory of a city delivered from the plague
at the prayer of its bishop, and of the Archangel who put up the sword of
God's judgment.

\paragraph{Mont-Saint-Michel.} The \emph{Golden Legend} counts the
appearing on the Norman coast as the second apparition of the Archangel:
``The second apparition was in the year of our Lord seven hundred and
ten, in a place which was named Tumba, by the seaside, six miles from the
city of Apricens. S.\ Michael appeared to the bishop of that city and
commanded him to do make a church in the foresaid place, like as it was
made in the mount of Gargan, and in like wise should hallow the memory of
S.\ Michael there'' (\emph{Golden Legend}, V, pp.~182--183). The bishop
of Avranches found the site where the Archangel showed him ``a bull hid
of thieves,'' and its bounds where the bull had trodden; a man sent by
Michael removed two rocks ``as lightly as they had weighed nothing''; and
the Archangel brought to the new church a piece of the marble on which he
had stood and part of the pall laid on the altar of Gargano. The \emph{Golden Legend}
adds that the apparition ``is solemnly hallowed the seventeenth kalends of
November in that place,'' that is, on 16 October, and tells of the woman
with child overtaken by the tide on the way to the church, whom Michael
kept whole, so that ``she was delivered and childed among the waves in the
middle of the sea'' (ibid., pp.~183--184). The Directory of 2001 names the
Mount first among the great shrines of the angels (\emph{Directory} 216).

\paragraph{The East.} The East commemorates an apparition of the
Archangel at Chonae in Phrygia on 6 September and keeps its solemnity of
the angels, the Synaxis of Saint Michael and the other heavenly powers, on
8 November (\emph{The Liturgical Year}, Time after Pentecost V, September
29; \S\ref{sec:byzantine}).

\sectionguard
\subsection{The Church's Rule for the Cult of the Angels}

The Church honours the angels with veneration and invocation, and she
directs all their honour to God. From the first centuries she has also
held the cult of the angels to what revelation gives: no worship apart
from the Church's worship, no names but the names of Scripture, and no
doctrine of the angels built on private revelations.

\paragraph{Laodicea.} The Apostle had warned the Colossians: ``Let no man
seduce you, willing in humility and religion of angels, walking in the
things which he hath not seen, in vain puffed up by the sense of his
flesh'' (Col 2:18), and bade them send his letter on to ``the church of
the Laodiceans'' (Col 4:16). In Laodicea a synod of the fourth century
forbade a cult of angels held apart from the Church:
\begin{quote}
Christians must not forsake the Church of God, and go away and invoke
angels and gather assemblies, which things are forbidden. If, therefore,
any one shall be found engaged in this covert idolatry, let him be
anathema; for he has forsaken our Lord Jesus Christ, the Son of God, and
has gone over to idolatry. (Laodicea, can.~35, trans.\ Percival, NPNF
II.14)
\end{quote}
What the canon condemns is the invocation of angels in assemblies that
forsake the Church and Christ; the Church's own invocation of the angels
at her altar is the opposite of that act. The Carolingian Church read the
canon as a rule for the names of the angels. Charlemagne's
\emph{Admonitio generalis} of 789, citing ``the same council'' of
Laodicea, directs \latin{ut ignota angelorum nomina nec fingantur nec
nominentur, nisi illos quos habemus in auctoritate, id sunt Michahel,
Gabrihel, Rafahel}: that unknown names of angels be neither invented nor
used, except those we have on authority, Michael, Gabriel, and Raphael
(\emph{Admonitio generalis}, c.~16, MGH \emph{Capit.}\ I, ed.\ Boretius, 1883, p.~55).

\paragraph{The Roman synod of 745.} The rule on names was first laid
down by a Roman synod. On 25 October 745 Pope Zachary held a synod in the
Lateran on the case of Aldebert, a Frankish cleric whom a synod under
Boniface had deprived of the priesthood and who was still leading the
people astray. At a later session Boniface's legate, the priest Deneard,
produced a prayer Aldebert had composed. It was read
until it reached the words, \latin{Praecor vos et coniuro vos et supplico
me ad vos, angelus Uriel, angelus Raguel, angelus Tubuel, angelus
Michael, angelus Adinus, angelus Tubuas, angelus Sabaoc, angelus Simiel}
(MGH \emph{Conc.}\ II/1, ed.\ Werminghoff, 1906, p.~42). The bishops and
priests answered that everything read should be burned and its authors
anathematized:
\begin{quote}
\latin{Quia octo nomina angelorum, quae in sua oratione Aldebertus
invocavit, non angelorum, praeterquam Michaelis, sed magis demones in sua
oratione sibi ad prestandum auxilium invocavit. Nos autem, ut a vestro
sancto apostolatu edocemur et divina tradit auctoritas, non plus quam
trium angelorum nomina cognoscimus, id est Michael, Gabriel, Raphael.}
(ibid., p.~43)
\end{quote}
Of the eight names Aldebert invoked, the synod judged, only Michael's is
an angel's; with the others he called on demons to help him. As the
Apostolic See teaches and divine authority hands down, the Church knows
the names of no more than three angels, Michael, Gabriel, and Raphael.
Zachary directed that the writings be kept in the Roman archive to
Aldebert's perpetual confusion, and the synod removed Aldebert from every
priestly office, to do penance, under anathema if he persisted (ibid.). The synod taught the rule with its reason:
names of angels not given by revelation open the door to the invocation of
demons.

\paragraph{The Directory on Popular Piety (2001).} The Congregation for
Divine Worship and the Discipline of the Sacraments restated the Church's
practice in its Directory of 17 December 2001. It grounds the devotion in
the faith and in the liturgy (213--215), describes its forms (216), and
names two deviations. The first is a view of the world ``subject to
demiurgical struggles, or an incessant battle between good and evil
spirits, or Angels and daemons, in which man is left at the mercy of
superior forces and over which he is helpless''; such cosmologies ``bear
little relation to the true Gospel vision of the struggle to overcome the
Devil, which requires moral commitment, a fundamental option for the
Gospel, humility and prayer.'' The second is a reading of daily life
``so as to ascribe all setbacks to the Devil and all success to the
Guardian Angels.'' It then states the rule on names: ``The practice of
assigning names to the Holy Angels should be discouraged, except in the
cases of Gabriel, Raphael and Michael whose names are contained in Holy
Scripture'' (\emph{Directory} 217).

\paragraph{The decree on the Opus Angelorum (1992).} The Congregation for
the Doctrine of the Faith applied the same rule to an association of the
faithful devoted to the angels. On 24 September 1983, with decisions
approved by the Pope on 1 July, it had directed that the association
follow the doctrine of the Church, the Fathers, and the Doctors in
promoting devotion to the angels, and not spread among its members a cult
of the angels using ``names'' known from the presumed private revelation
attributed to Gabriele Bitterlich, nor use those names in any prayer
(\emph{Decretum de doctrina et usibus particularibus consociationis cui
nomen ``Opus Angelorum''}, 6 June 1992, recalling AAS 76 [1984]
175--176). Having examined further writings, the Congregation judged that
the association's own angelology and some usages derived from it were
\latin{alienos esse a S.\ Scriptura necnon a Traditione}, foreign to
Scripture and Tradition, and could not ground the spirituality of
associations approved by the Church. With the approval of John Paul~II it
decreed:
\begin{quote}
\latin{Theoriae provenientes ex praesumptis revelationibus Dominae
Gabrielae Bitterlich de mundo angelorum, de eorum nominibus personalibus,
coetibus et functionibus non possunt nec doceri nec ullo modo adhiberi,
sive explicite sive implicite~\dots\ in cultu, in orationibus, in
formatione spirituali, in spiritualitate publica vel privata, in
ministerio et apostolatu.} (ibid., I; AAS 84 [1992] 805--806)
\end{quote}
Theories about the world of the angels, their personal names, their
groupings, and their functions, drawn from these presumed revelations,
may be neither taught nor used in any way. The decree forbade the
consecrations to angels practised in the association, the so-called
administration of the sacraments at a distance, and the insertion into
the Mass and the Office of texts or rites referring to those theories; it
directed that exorcisms be used only according to the norms of the
Church and with the formulas she approves; and it named a delegate of the
Holy See to see to the decree's execution (ibid., II--V). For the Church's
teaching on the cult of the angels the decree cites Benedict~XIV's
treatise on the beatification and canonization of the servants of God,
book IV, part 2, chapter 30, \latin{De Angelis et eorum cultu} (ibid.,
n.~1).

\paragraph{The rule.} The acts agree. The Church venerates the angels
and invokes them, in the liturgy and in private prayer, because God has
sent them to minister to those who shall receive the inheritance of
salvation (Heb 1:14). She venerates them in her own worship, which is
the worship of God through Christ: the Preface praises the Father
``Thro' Christ our Lord: by whom the angels praise thy majesty,'' and the
prayer to Michael asks him to act ``by the divine power.'' She names the
angels whom Scripture names, Michael, Gabriel, and Raphael, invokes the
rest by their orders, and forbids the invention of other names, which the
Roman synod of 745 judged a door to the invocation of demons. She does not
build her doctrine of the angels on private revelations, and she holds
the devotion within the Gospel's account of the struggle with the devil,
which is won by faith, humility, and prayer.

\Needspace{14\baselineskip}
\subsection{Devotion and Its Rule, by Degree}

\begin{threestable}{Grade}{Proposition}{Basis}
Church teaching & The Church venerates the angels and has recourse to
their intercession; the faithful may invoke them in private prayer &
CCC~335; \emph{Directory} 215--216; the indulgenced prayers of the
\emph{Raccolta}\\
Church teaching & Each of the faithful has a guardian angel, to be
honoured with reverence, gratitude, and confidence & CCC~336;
\emph{Directory} 216; collect of 2 Oct.; Bernard, lesson v of
2 Oct.\\
Church teaching & A cult of angels apart from the Church's worship is
forbidden & Laodicea, can.~35 (disciplinary canon of a particular
synod)\\
Church teaching & Only the names Scripture gives, Michael, Gabriel,
Raphael, are to be used; other names are not to be invented or invoked &
Roman synod of 745 (MGH \emph{Conc.}\ II/1, p.~43); \emph{Admonitio
generalis} c.~16 (789); \emph{Directory} 217 (disciplinary acts)\\
Church teaching & Doctrines on angels' names, choirs, and functions drawn
from private revelations may not be taught or used; consecrations to
angels in the Opus Angelorum are forbidden & CDF decree of 6 June 1992
(AAS 84 [1992] 805--806; disciplinary decree)\\
Common teaching & Michael is the prince and defender of the Church and
the Church's champion against the devil & Dan 10:21, 12:1; Apoc 12:7--9;
Breviary of 29 Sept.; prayer after Low Mass\\
Pious tradition & The Archangel appeared on Mount Gargano under
Gelasius~I and asked that God be worshipped in the cave & Breviary,
lessons of 8 May; Martyrology of 29 Sept.; \emph{Golden Legend}\\
Pious tradition & Gregory the Great saw the angel sheathe his sword over
Hadrian's mausoleum at the end of the plague & \emph{Golden Legend},
St Michael\\
Pious tradition & The Archangel appeared to the bishop of Avranches and
asked for a church at Tumba, Mont-Saint-Michel & \emph{Golden Legend},
St Michael; \emph{Directory} 216\\
\end{threestable}
